\section{总结、当前边界与后续演进}

\subsection{核心总结}

回到第 1 章提出的自由行规划四大痛点，本项目围绕“\textbf{可核对、可微调、本地可控、易安装维护}”四项原则，完成了从需求收敛、架构决策、核心引擎研发到单容器交付的完整工程实现：

\begin{itemize}
  \item \textbf{架构与工程扎实}：以单应用容器统一交付 React + TypeScript 前端与 FastAPI + SQLite (FTS5, Schema v6) 后端，配套进程独占锁、事务迁移与 \texttt{travel-data} 离线校验/备份/空卷恢复工具，同时支持 Docker Compose 一键部署与 \texttt{git clone} 本机运行。
  \item \textbf{规划与核验可信}：动态 LLM 动作循环配合程序规则校验、独立质量审核与检查点续跑，彻底告别了“一错全废”或“盲目信任大模型”的两难；分离式三段原文引用与 \texttt{Exa $\rightarrow$ DeepSeek $\rightarrow$ Tavily $\rightarrow$ AnySearch} 四家搜索回退，使营业时间、淡旺季票价与预约要求真正做到有据可查、未知不冒充通过。
  \item \textbf{体验克制且便于落地}：自由用餐时段、纯净活动介绍、定点局部编辑与只读历史版本让行程真正具备日常可操作性；单人本机存储、自然语言长期偏好与行程级按需资料关联，在个性化效果与个人隐私边界之间取得了清晰平衡。
\end{itemize}

\subsection{当前已知边界}

受限于单机架构取舍与外部公开信息生态，当前版本仍存在四类客观边界，适合在使用与后续迭代中持续关注：

\begin{enumerate}
  \item \textbf{多轮串行规划耗时相对较长}：在外部模型串行决策、联网取证与独立审核模式下，多日复杂行程涉及较多外部调用轮次，整体生成耗时相对较长，偶发单步请求超时（系统通过检查点续跑与降级机制保底，长行程建议适当调大 \texttt{PROVIDER\_TIMEOUT\_SECONDS}，例如设为 120 秒）。
  \item \textbf{官方事实覆盖仍受公开搜索生态制约}：虽然部分大型核心场馆能够核实官方淡季票价与预约要求，但公开搜索引擎返回结果仍以第三方攻略和聚合页居多，导致行程中仍有相当比例的营业时间与费用停留在 \texttt{estimated} 或 \texttt{unknown}，最终呈现为 \texttt{needs\_review}。
  \item \textbf{跨市区与远郊排程的自修复空间}：当行程同时涵盖市中心景点与远郊景区且叠加较严格的步行距离偏好时，模型偶发仍可能生成不合理的活动间隔或交通方式；虽然程序规则校验会可靠拦截为 \texttt{blocked} 避免误导用户，但草案在修复轮次中的自动纠偏能力仍有进一步提升空间。
  \item \textbf{交互与功能边界留白}：当前偏好更新采用用户在设置页主动编辑保存的方式，尚未在行程编辑后自动弹出“建议写入长期偏好”的确认条；天气预报、多城串联行程以及扫描件 PDF 的本地 OCR 均按首版简化原则留待后续迭代。
\end{enumerate}

\subsection{后续演进方向}

面向下一阶段的实用化打磨，建议按投入产出比依次推进以下四项工作：

\begin{itemize}
  \item \textbf{取证并发与官方站点定向召回}：在遵守各搜索提供方 QPS 与冷却机制的前提下，对独立场馆的事实查询引入受控并发并强化对官网/政府域名（如 \texttt{.gov.cn}、景区官网）的优先检索策略，同时缩短整体规划时长并提升 \texttt{supported} 事实占比。
  \item \textbf{远郊与跨区交通规则预检}：在草案生成后、消耗修复轮次前增加显式的长距离步行纠偏提示（例如超过 3 公里自动回传公交/驾车耗时对照给模型），减少远郊跨区场景下的交通冲突。
  \item \textbf{偏好提炼确认交互}：在局部编辑完成后可选异步提炼一条简短的长期偏好更新建议，以非侵入式提示条展示在页面，由用户一键确认或忽略，进一步打通 R03 的体验闭环。
  \item \textbf{可选天气与使用反馈闭环}：按既有接口预留接入和风天气 7--10 天预报作为室内外排程参考，并在日常出行规划中持续以实际使用反馈打磨规则与提示词。
\end{itemize}
